% Options for packages loaded elsewhere
% Options for packages loaded elsewhere
\PassOptionsToPackage{unicode}{hyperref}
\PassOptionsToPackage{hyphens}{url}
\PassOptionsToPackage{dvipsnames,svgnames,x11names}{xcolor}
%
\documentclass[
  authoryear,
  preprint,
  3p]{elsarticle}
\usepackage{xcolor}
\usepackage{amsmath,amssymb}
\setcounter{secnumdepth}{5}
\usepackage{iftex}
\ifPDFTeX
  \usepackage[T1]{fontenc}
  \usepackage[utf8]{inputenc}
  \usepackage{textcomp} % provide euro and other symbols
\else % if luatex or xetex
  \usepackage{unicode-math} % this also loads fontspec
  \defaultfontfeatures{Scale=MatchLowercase}
  \defaultfontfeatures[\rmfamily]{Ligatures=TeX,Scale=1}
\fi
\usepackage{lmodern}
\ifPDFTeX\else
  % xetex/luatex font selection
\fi
% Use upquote if available, for straight quotes in verbatim environments
\IfFileExists{upquote.sty}{\usepackage{upquote}}{}
\IfFileExists{microtype.sty}{% use microtype if available
  \usepackage[]{microtype}
  \UseMicrotypeSet[protrusion]{basicmath} % disable protrusion for tt fonts
}{}
\makeatletter
\@ifundefined{KOMAClassName}{% if non-KOMA class
  \IfFileExists{parskip.sty}{%
    \usepackage{parskip}
  }{% else
    \setlength{\parindent}{0pt}
    \setlength{\parskip}{6pt plus 2pt minus 1pt}}
}{% if KOMA class
  \KOMAoptions{parskip=half}}
\makeatother
% Make \paragraph and \subparagraph free-standing
\makeatletter
\ifx\paragraph\undefined\else
  \let\oldparagraph\paragraph
  \renewcommand{\paragraph}{
    \@ifstar
      \xxxParagraphStar
      \xxxParagraphNoStar
  }
  \newcommand{\xxxParagraphStar}[1]{\oldparagraph*{#1}\mbox{}}
  \newcommand{\xxxParagraphNoStar}[1]{\oldparagraph{#1}\mbox{}}
\fi
\ifx\subparagraph\undefined\else
  \let\oldsubparagraph\subparagraph
  \renewcommand{\subparagraph}{
    \@ifstar
      \xxxSubParagraphStar
      \xxxSubParagraphNoStar
  }
  \newcommand{\xxxSubParagraphStar}[1]{\oldsubparagraph*{#1}\mbox{}}
  \newcommand{\xxxSubParagraphNoStar}[1]{\oldsubparagraph{#1}\mbox{}}
\fi
\makeatother


\usepackage{longtable,booktabs,array}
\usepackage{calc} % for calculating minipage widths
% Correct order of tables after \paragraph or \subparagraph
\usepackage{etoolbox}
\makeatletter
\patchcmd\longtable{\par}{\if@noskipsec\mbox{}\fi\par}{}{}
\makeatother
% Allow footnotes in longtable head/foot
\IfFileExists{footnotehyper.sty}{\usepackage{footnotehyper}}{\usepackage{footnote}}
\makesavenoteenv{longtable}
\usepackage{graphicx}
\makeatletter
\newsavebox\pandoc@box
\newcommand*\pandocbounded[1]{% scales image to fit in text height/width
  \sbox\pandoc@box{#1}%
  \Gscale@div\@tempa{\textheight}{\dimexpr\ht\pandoc@box+\dp\pandoc@box\relax}%
  \Gscale@div\@tempb{\linewidth}{\wd\pandoc@box}%
  \ifdim\@tempb\p@<\@tempa\p@\let\@tempa\@tempb\fi% select the smaller of both
  \ifdim\@tempa\p@<\p@\scalebox{\@tempa}{\usebox\pandoc@box}%
  \else\usebox{\pandoc@box}%
  \fi%
}
% Set default figure placement to htbp
\def\fps@figure{htbp}
\makeatother





\setlength{\emergencystretch}{3em} % prevent overfull lines

\providecommand{\tightlist}{%
  \setlength{\itemsep}{0pt}\setlength{\parskip}{0pt}}



 
\usepackage[]{natbib}


\usepackage{booktabs}
\usepackage{longtable}
\usepackage{array}
\usepackage{multirow}
\usepackage{wrapfig}
\usepackage{float}
\usepackage{colortbl}
\usepackage{pdflscape}
\usepackage{tabu}
\usepackage{threeparttable}
\usepackage{threeparttablex}
\usepackage[normalem]{ulem}
\usepackage{makecell}
\usepackage{xcolor}
\usepackage{pdflscape}
\makeatletter
\@ifpackageloaded{caption}{}{\usepackage{caption}}
\AtBeginDocument{%
\ifdefined\contentsname
  \renewcommand*\contentsname{Table of contents}
\else
  \newcommand\contentsname{Table of contents}
\fi
\ifdefined\listfigurename
  \renewcommand*\listfigurename{List of Figures}
\else
  \newcommand\listfigurename{List of Figures}
\fi
\ifdefined\listtablename
  \renewcommand*\listtablename{List of Tables}
\else
  \newcommand\listtablename{List of Tables}
\fi
\ifdefined\figurename
  \renewcommand*\figurename{Figure}
\else
  \newcommand\figurename{Figure}
\fi
\ifdefined\tablename
  \renewcommand*\tablename{Table}
\else
  \newcommand\tablename{Table}
\fi
}
\@ifpackageloaded{float}{}{\usepackage{float}}
\floatstyle{ruled}
\@ifundefined{c@chapter}{\newfloat{codelisting}{h}{lop}}{\newfloat{codelisting}{h}{lop}[chapter]}
\floatname{codelisting}{Listing}
\newcommand*\listoflistings{\listof{codelisting}{List of Listings}}
\makeatother
\makeatletter
\usepackage{pdflscape}
\makeatother
\makeatletter
\makeatother
\makeatletter
\@ifpackageloaded{caption}{}{\usepackage{caption}}
\@ifpackageloaded{subcaption}{}{\usepackage{subcaption}}
\makeatother
\journal{Transfiguration Today}
\usepackage{bookmark}
\IfFileExists{xurl.sty}{\usepackage{xurl}}{} % add URL line breaks if available
\urlstyle{same}
\hypersetup{
  pdftitle={An Empirical Study of Animagus Transformations and Their Metabolic Signatures},
  pdfauthor={Harry Potter; Hermione Granger},
  pdfkeywords={Transfiguration, Animagus, Hogwarts, Magic},
  colorlinks=true,
  linkcolor={blue},
  filecolor={Maroon},
  citecolor={Blue},
  urlcolor={Blue},
  pdfcreator={LaTeX via pandoc}}


\setlength{\parindent}{6pt}
\begin{document}

\begin{frontmatter}
\title{An Empirical Study of Animagus Transformations and Their Metabolic Signatures}
\author[1,2]{Harry Potter%
\corref{cor1}%
}
 \ead{harrypotter@hogwarts.edu} 
\author[2]{Hermione Granger%
%
}
 \ead{hermione.granger@ministryofmagic.gov.uk} 

\affiliation[1]{organization={Hogwarts School of Witchcraft and Wizardry, Gryffindor House},city={Highlands},country={United Kingdom},countrysep={,},postcodesep={}}
\affiliation[2]{organization={Ministry of Magic, Department of Magical Law Enforcement},city={London},country={United Kingdom},countrysep={,},postcodesep={}}

\cortext[cor1]{Corresponding author}


        
\begin{abstract}
Animagus transformation represents one of the most advanced branches of Transfiguration. While the qualitative protocols for achieving such transformations are well-documented, the underlying metabolic pathways and physiological signatures remain poorly understood. In this study, we monitor the metabolic activity of seasoned Animagi during the transition between human and animal forms. Our findings demonstrate a temporary but significant spike in magical energy consumption and local molecular transmutations, which are tightly regulated by standard transfiguration equations. This research provides a novel quantitative foundation for the biophysics of magical shapeshifting.
\end{abstract}





\begin{keyword}
    Transfiguration \sep Animagus \sep Hogwarts \sep 
    Magic
\end{keyword}
\end{frontmatter}
    

\section{Introduction}\label{introduction}

This repository provides a highly automated, programmatic workflow for academic writing via Quarto \citep{smith2023}.
It replaces the traditional ``copy and paste'' approach with dynamic data embedding, guaranteeing strict reproducibility and minimizing human error when migrating analytical results from external scripts into the final manuscript.

To achieve robust modularity, we employ a central data bridging directory (configured via \texttt{user-settings} in \texttt{\_quarto.yml}, defaulting to \texttt{sample\_data/}). Analysis pipelines (e.g., Python or R scripts) can blindly export their payloads into this folder, and the Quarto compiler will symmetrically bind and format them during document rendering.

\section{Methods}\label{methods}

\subsection{Dynamic Data Embedding}\label{dynamic-data-embedding}

A foundational feature of this framework is the automated resolution of scalar outputs. Rather than hardcoding data points into the document body, external YAML dictionaries (\texttt{results.yaml}) are statically mounted into the native \texttt{R} environment upon compilation.

Custom YAML tags (e.g., \texttt{!iqFloat}, \texttt{!iqPval}) paired with vectorized type-matching handlers guarantee that results are rigorously rounded, stripped of leading zeros, and formulated to Elsevier manuscript standards without requiring sprawling inline parameter constraints.

\subsection{Variable Injection Demo}\label{variable-injection-demo}

As a practical example, we can seamlessly print explicit statistics into the running text using R inline blocks. In our mock evaluation, the total sample size was \(N = 1450\), and the mean age was \(35.4\) years.

Our primary analysis demonstrated an effect size of \(d = 0.85\) (\(p = .003\)) with a sustained \(R^2\) of \(0.7329488\). We further verified granular correlations across the matrix: \(r = .492\) (\(p < .001\)) and \(r = -.322\). All variables were imported asynchronously and bound via \texttt{tools/import\_r\_tools.R}.

\clearpage

\section{Tables}\label{tables}

\clearpage

\begin{landscape}

\begin{table}

\caption{\label{tbl-example}\textbf{Automated MultiIndex Tables.}\\
\strut \\
This table compiles directly from \texttt{sample\_data/table/table\_data.csv}. The backend R scripts auto-detect CSV arrays dropped into the \texttt{table\_dir}, mapping column string annotations (\texttt{::iqIndex}, \texttt{::iqFloat}) to natively construct nested booktabs without manual row-collapse logic. Note how \(p\)-values and mathematical strings like \(\chi^2\) are cleanly resolved horizontally through the \texttt{make\_table()} utility.}

\centering{

\begingroup\renewcommand{\arraystretch}{1.15}
\setlength{\defaultaddspace}{1.0em}

\begin{tabular}[t]{lllllll}
\toprule
 &  & mean\_val & sd\_val & \textit{r}\textsubscript{Pearson} & p\_corr & $\chi^2$\\
\midrule
 & Baseline & 18.24 & 1.218392 & .124 & .346 & 15.23\\

\multirow[t]{-2}{*}[1\dimexpr\aboverulesep+\belowrulesep+\cmidrulewidth]{\raggedright\arraybackslash Control} & treat & 22.19 & 1.529384 & .234 & .013 & 18.99\\
\addlinespace
 & Baseline & 24.32 & 1.829103 & -.052 & .884 & 8.12\\

\multirow[t]{-2}{*}[1\dimexpr\aboverulesep+\belowrulesep+\cmidrulewidth]{\raggedright\arraybackslash Group A} & treat & 31.53 & 2.193812 & .492 & < .001 & 2.44\\
\addlinespace
 & Baseline & 12.02 & 0.928193 & -.322 & .004 & 4.19\\

\multirow[t]{-2}{*}[1\dimexpr\aboverulesep+\belowrulesep+\cmidrulewidth]{\raggedright\arraybackslash Group B} & treat & 29.19 & 2.019283 & .850 & < .001 & 26.54\\
\bottomrule
\end{tabular}
\endgroup

}

\end{table}%

\clearpage

\end{landscape}

\clearpage

\section{Figures}\label{figures}

\begin{figure}

\centering{

\pandocbounded{\includegraphics[keepaspectratio]{mock_plot.png}}

}

\caption{\label{fig-example}\textbf{Native Image Handling Pipeline.}\\
\strut \\
When an external figure is exported into the \texttt{figure\_dir} folder (\texttt{sample\_data/figure}), Quarto handles the indexing, but the Elsevier LaTeX compiler notoriously struggles with nested relative paths. Our \texttt{tools/copy\_figures.R} pre-flight hook actively solves this by transferring requested assets momentarily toward the root compilation context.}

\end{figure}%

\clearpage

\section{Conclusion}\label{conclusion}

By isolating complex metadata extraction, tabular configuration, and figure management into the \texttt{tools/} backend loop, the user writing the actual \texttt{manuscript.qmd} is free to use pure Markdown functionality, guaranteeing publication-ready PDFs equipped with CRediT contribution trackers at every render sequence.

\section{Acknowledgements}\label{acknowledgements}

The authors wish to thank the Quarto development community and the R Foundation for Statistical Computing for creating robust scholarly scaffolding tools.

\clearpage

\section{Author Contributions}\label{author-contributions}

\textbf{Harry Potter}: Conceptualization, Methodology, Writing - Original Draft.

\textbf{Hermione Granger}: Supervision, Writing - Review \& Editing.

\clearpage

\section{References}\label{references}

\renewcommand{\bibsection}{}
\begin{thebibliography}{1}
\expandafter\ifx\csname natexlab\endcsname\relax\def\natexlab#1{#1}\fi
\providecommand{\url}[1]{\texttt{#1}}
\providecommand{\href}[2]{#2}
\providecommand{\path}[1]{#1}
\providecommand{\DOIprefix}{doi:}
\providecommand{\ArXivprefix}{arXiv:}
\providecommand{\URLprefix}{URL: }
\providecommand{\Pubmedprefix}{pmid:}
\providecommand{\doi}[1]{\href{http://dx.doi.org/#1}{\path{#1}}}
\providecommand{\Pubmed}[1]{\href{pmid:#1}{\path{#1}}}
\providecommand{\bibinfo}[2]{#2}
\ifx\xfnm\relax \def\xfnm[#1]{\unskip,\space#1}\fi
%Type = Article
\bibitem[{Smith and Doe(2023)}]{smith2023}
\bibinfo{author}{Smith, J.}, \bibinfo{author}{Doe, J.}, \bibinfo{year}{2023}.
\newblock \bibinfo{title}{A lorem ipsum analysis of quantitative methodologies}.
\newblock \bibinfo{journal}{Journal of Quantitative Lorem Ipsum} \bibinfo{volume}{42}, \bibinfo{pages}{1--15}.

\end{thebibliography}





\end{document}
